\documentclass[12pt, ukrainian]{article}
\usepackage[a4paper]{geometry}
\setlength\parindent{0mm}
%\setlength\parskip{4mm}
\geometry{top=1.3cm,bottom=1.3cm,left=1.5cm,right=1.5cm}
\pagestyle{empty}
\usepackage[T2A]{fontenc}
\usepackage[utf8]{inputenc}
\usepackage[ukrainian]{babel}
\begin{document}
%begin
{{\begin {scriptsize}\par
{ \center  Київський національний університет імені Тараса Шевченка \par}
{\parbox {5 cm}{Освітня програма \hfill}{\parbox {7 cm}{Прикладна математика\hfill}}\parbox {6 cm}{\hfill Семестр 2}}\par
{\parbox {5 cm} {Навчальний предмет \hfill}\parbox {7 cm}{Програмування \hfill}\parbox {6cm}{\hfill}\par}\vspace{-15pt} \end{scriptsize}}{\center Екзаменаційний білет \No \textbf { 1 }\par}}
{%beginitem
1. Багатовимірні масиви бібліотеки numpy: у чому відмінність моделювання матриці багатовимірним масивом бібліотеки numpy та як список списків.%enditem
\par
%beginitem
2. Написати функцію, що за текстовим файлом з статтею визначає  скільки разів у ньому містяться слова, що закінчуються на \textbf{ський} або на \textbf{зький} (рід, число та відмінок значення не мають, тобто \textit{міська}, \textit{міському} тощо теж мають рахуватись). Файл дуже великий і його вміст завантажений у пам'ять бути не може.

%enditem
\par
%beginitem
3. Реалізувати операцію злиття множин з повтореннями. Множина з повтореннями задається переліком своїх елементів, записаних за неспаданням у текстовому файлі через білі символи. Результат у такому ж самому форматі (і теж за неспаданням) має бути записаний у текстовий файл. Питома функція має отримувати три шляхи до файлів. Розмір вхідних множин такий, що вони не можуть бути завантажені повністю в пам’ять програми.

%enditem
}
{\smallskip\smallskip \begin {scriptsize} Затверджено на засіданні кафедри теоретичної кібернетики, протокол \No 4  від   25.01.2024 р.\par\smallskip\smallskip\smallskip\smallskip
{\parbox {6.5 cm} {Зав. кафедри \dotfill Ю.В.~Крак}}\hspace {0.7cm} {\hbox to 6.5 cm { Екзаменатор \dotfill Т.О.~Карнаух}} \end{scriptsize}}
%end
\vspace{0.9cm}
%begin
{{\begin {scriptsize}\par
{ \center  Київський національний університет імені Тараса Шевченка \par}
{\parbox {5 cm}{Освітня програма \hfill}{\parbox {7 cm}{Прикладна математика\hfill}}\parbox {6 cm}{\hfill Семестр 2}}\par
{\parbox {5 cm} {Навчальний предмет \hfill}\parbox {7 cm}{Програмування \hfill}\parbox {6cm}{\hfill}\par}\vspace{-15pt} \end{scriptsize}}{\center Екзаменаційний білет \No \textbf { 2 }\par}}
{%beginitem
1. Алгоритм пошуку в глибину та задачі, які він може вирішувати.%enditem
\par
%beginitem
2. Написати функцію, що в текстовому файлі замінює всі цілі числа на слово NUMBER. Файл дуже великий і його вміст завантажений у пам'ять бути не може. Результат має бути записаний у новий файл.

%enditem
\par
%beginitem
3. Реалізувати операцію різниці множин з повтореннями. Множина з повтореннями задається переліком своїх елементів, записаних за неспаданням у текстовому файлі через білі символи. Результат у такому ж самому форматі (і теж за неспаданням) має бути записаний у текстовий файл. Питома функція має отримувати три шляхи до файлів. Розмір вхідних множин такий, що вони не можуть бути завантажені повністю в пам’ять програми.

%enditem
}
{\smallskip\smallskip \begin {scriptsize} Затверджено на засіданні кафедри теоретичної кібернетики, протокол \No 4  від   25.01.2024 р.\par\smallskip\smallskip\smallskip\smallskip
{\parbox {6.5 cm} {Зав. кафедри \dotfill Ю.В.~Крак}}\hspace {0.7cm} {\hbox to 6.5 cm { Екзаменатор \dotfill Т.О.~Карнаух}} \end{scriptsize}}
%end
\vspace{0.9cm}
%begin
{{\begin {scriptsize}\par
{ \center  Київський національний університет імені Тараса Шевченка \par}
{\parbox {5 cm}{Освітня програма \hfill}{\parbox {7 cm}{Прикладна математика\hfill}}\parbox {6 cm}{\hfill Семестр 2}}\par
{\parbox {5 cm} {Навчальний предмет \hfill}\parbox {7 cm}{Програмування \hfill}\parbox {6cm}{\hfill}\par}\vspace{-15pt} \end{scriptsize}}{\center Екзаменаційний білет \No \textbf { 3 }\par}}
{%beginitem
1. Файли формату json.%enditem
\par
%beginitem
2. Написати функцію, що в текстовому файлі замінює всі зазначення ваги у вигляді \textbf{якесь ціле число кг} на слово WEIGHT (приклади зазначення ваги: 13 кг, 26   кг). Файл дуже великий і його вміст завантажений у пам'ять бути не може. Результат має бути записаний у новий файл.

%enditem
\par
%beginitem
3. Текстовий файл містить послідовність цілих чисел $a_1$, $a_2$, $\ldots$, $a_t$, записану у форматі одне число на рядок. Знайти сімнадцяте найбільше значення, що записане у файлі. Кожне значення, що записане у файлі, враховується стільки разів, скільки входить у файл. Кількість чисел у файлі така, що вони не можуть бути завантажені повністю в пам'ять програми. Розв'язання оформити у вигляді функції.

%enditem
}
{\smallskip\smallskip \begin {scriptsize} Затверджено на засіданні кафедри теоретичної кібернетики, протокол \No 4  від   25.01.2024 р.\par\smallskip\smallskip\smallskip\smallskip
{\parbox {6.5 cm} {Зав. кафедри \dotfill Ю.В.~Крак}}\hspace {0.7cm} {\hbox to 6.5 cm { Екзаменатор \dotfill Т.О.~Карнаух}} \end{scriptsize}}
%end
\newpage
%begin
{{\begin {scriptsize}\par
{ \center  Київський національний університет імені Тараса Шевченка \par}
{\parbox {5 cm}{Освітня програма \hfill}{\parbox {7 cm}{Прикладна математика\hfill}}\parbox {6 cm}{\hfill Семестр 2}}\par
{\parbox {5 cm} {Навчальний предмет \hfill}\parbox {7 cm}{Програмування \hfill}\parbox {6cm}{\hfill}\par}\vspace{-15pt} \end{scriptsize}}{\center Екзаменаційний білет \No \textbf { 4 }\par}}
{%beginitem
1. Поняття стеку та черги. Основні операції зі стеками та чергами.%enditem
\par
%beginitem
2. Написати функцію, що за текстовим файлом, що містить зображення цілих чисел, дійсних чисел та точок площини, розділені білими символами, обчислює кількість входжень у файл зображень точок. Зображення точки має вигляд: відкриваюча дужка, x-координата, кома, y-координата, закриваюча дужка, наприклад, ( 13.5 , 45.1E-1) або ( 14, -3) . Білі символи можуть зустрічатись у довільній кількості. Файл дуже великий і його вміст завантажений у пам'ять бути не може.

%enditem
\par
%beginitem
3. Реалізувати операцію симетричної різниці множин з повтореннями. Множина з повтореннями задається переліком своїх елементів, записаних за неспаданням у текстовому файлі через білі символи. Результат у такому ж самому форматі (і теж за неспаданням) має бути записаний у текстовий файл. Питома функція має отримувати три шляхи до файлів. Розмір вхідних множин такий, що вони не можуть бути завантажені повністю в пам’ять програми.

%enditem
}
{\smallskip\smallskip \begin {scriptsize} Затверджено на засіданні кафедри теоретичної кібернетики, протокол \No 4  від   25.01.2024 р.\par\smallskip\smallskip\smallskip\smallskip
{\parbox {6.5 cm} {Зав. кафедри \dotfill Ю.В.~Крак}}\hspace {0.7cm} {\hbox to 6.5 cm { Екзаменатор \dotfill Т.О.~Карнаух}} \end{scriptsize}}
%end
\vspace{0.9cm}
%begin
{{\begin {scriptsize}\par
{ \center  Київський національний університет імені Тараса Шевченка \par}
{\parbox {5 cm}{Освітня програма \hfill}{\parbox {7 cm}{Прикладна математика\hfill}}\parbox {6 cm}{\hfill Семестр 2}}\par
{\parbox {5 cm} {Навчальний предмет \hfill}\parbox {7 cm}{Програмування \hfill}\parbox {6cm}{\hfill}\par}\vspace{-15pt} \end{scriptsize}}{\center Екзаменаційний білет \No \textbf { 5 }\par}}
{%beginitem
1. Алгоритм пошуку в ширину та задачі, які він може вирішувати.%enditem
\par
%beginitem
2. Написати функцію, що за текстовим файлом з статтею визначає скільки разів у ньому містяться слова, що починаються на Київ та Kyiv.  Переноси слів між рядками у файлі відсутні. Файл дуже великий і його вміст завантажений у пам'ять бути не може.

%enditem
\par
%beginitem
3. Текстовий файл містить послідовність цілих чисел $a_1$, $a_2$, $\ldots$, $a_t$, записану у форматі одне число на рядок. Знайти сімнадцяте найменше значення, що записане у файлі. Кожне значення, що записане у файлі, враховується стільки разів, скільки входить у файл. Кількість чисел у файлі така, що вони не можуть бути завантажені повністю в пам'ять програми. Розв'язання оформити у вигляді функції.

%enditem
}
{\smallskip\smallskip \begin {scriptsize} Затверджено на засіданні кафедри теоретичної кібернетики, протокол \No 4  від   25.01.2024 р.\par\smallskip\smallskip\smallskip\smallskip
{\parbox {6.5 cm} {Зав. кафедри \dotfill Ю.В.~Крак}}\hspace {0.7cm} {\hbox to 6.5 cm { Екзаменатор \dotfill Т.О.~Карнаух}} \end{scriptsize}}
%end
\vspace{0.9cm}
%begin
{{\begin {scriptsize}\par
{ \center  Київський національний університет імені Тараса Шевченка \par}
{\parbox {5 cm}{Освітня програма \hfill}{\parbox {7 cm}{Прикладна математика\hfill}}\parbox {6 cm}{\hfill Семестр 2}}\par
{\parbox {5 cm} {Навчальний предмет \hfill}\parbox {7 cm}{Програмування \hfill}\parbox {6cm}{\hfill}\par}\vspace{-15pt} \end{scriptsize}}{\center Екзаменаційний білет \No \textbf { 6 }\par}}
{%beginitem
1. Контекстні менеджери.%enditem
\par
%beginitem
2. Написати функцію, що в текстовому файлі замінює всі входження фрази \textbf{не зараховано} на слово UPS. Кількість білих символів між частинами фрази може бути довільна (і там можуть бути не тільки пробіли). Результат має бути записаний у новий файл.

%enditem
\par
%beginitem
3. Реалізувати операцію перетину множин з повтореннями. Множина з повтореннями задається переліком своїх елементів, записаних за неспаданням у текстовому файлі через білі символи. Результат у такому ж самому форматі (і теж за неспаданням) має бути записаний у текстовий файл. Питома функція має отримувати три шляхи до файлів. Розмір вхідних множин такий, що вони не можуть бути завантажені повністю в пам’ять програми.

%enditem
}
{\smallskip\smallskip \begin {scriptsize} Затверджено на засіданні кафедри теоретичної кібернетики, протокол \No 4  від   25.01.2024 р.\par\smallskip\smallskip\smallskip\smallskip
{\parbox {6.5 cm} {Зав. кафедри \dotfill Ю.В.~Крак}}\hspace {0.7cm} {\hbox to 6.5 cm { Екзаменатор \dotfill Т.О.~Карнаух}} \end{scriptsize}}
%end
\newpage
%begin
{{\begin {scriptsize}\par
{ \center  Київський національний університет імені Тараса Шевченка \par}
{\parbox {5 cm}{Освітня програма \hfill}{\parbox {7 cm}{Прикладна математика\hfill}}\parbox {6 cm}{\hfill Семестр 2}}\par
{\parbox {5 cm} {Навчальний предмет \hfill}\parbox {7 cm}{Програмування \hfill}\parbox {6cm}{\hfill}\par}\vspace{-15pt} \end{scriptsize}}{\center Екзаменаційний білет \No \textbf { 7 }\par}}
{%beginitem
1. Поняття часової та просторової складності обчислень. Складність у середньому. Асимптотичні оцінки складності.%enditem
\par
%beginitem
2. Написати функцію, що для текстового файлу, що містить послідовність точок площини, обчислює центр мас заданих точок (центр мас -- точка, координати якої є середнім арифметичним відповідних координат). Для кожної точки на окремому рядку наводиться її зображення: відкриваюча дужка, x-координата, кома, y-координата, закриваюча дужка, наприклад, ( 13.5 , 45.1E-1). Білі символи можуть зустрічатись у довільній кількості. Файл дуже великий і його вміст завантажений у пам'ять бути не може.

%enditem
\par
%beginitem
3. Кожен з двох текстових файлів містить послідовність дійсних чисел, записану у форматі одне число на рядок.  Перевірити, чи задають вони одну ту саму послідовність. Кількість чисел у кожному з двох файлів така, що вони не можуть бути завантажені повністю в пам'ять програми. Розв'язання оформити у вигляді функції.

%enditem
}
{\smallskip\smallskip \begin {scriptsize} Затверджено на засіданні кафедри теоретичної кібернетики, протокол \No 4  від   25.01.2024 р.\par\smallskip\smallskip\smallskip\smallskip
{\parbox {6.5 cm} {Зав. кафедри \dotfill Ю.В.~Крак}}\hspace {0.7cm} {\hbox to 6.5 cm { Екзаменатор \dotfill Т.О.~Карнаух}} \end{scriptsize}}
%end
\vspace{0.9cm}
%begin
{{\begin {scriptsize}\par
{ \center  Київський національний університет імені Тараса Шевченка \par}
{\parbox {5 cm}{Освітня програма \hfill}{\parbox {7 cm}{Прикладна математика\hfill}}\parbox {6 cm}{\hfill Семестр 2}}\par
{\parbox {5 cm} {Навчальний предмет \hfill}\parbox {7 cm}{Програмування \hfill}\parbox {6cm}{\hfill}\par}\vspace{-15pt} \end{scriptsize}}{\center Екзаменаційний білет \No \textbf { 8 }\par}}
{%beginitem
1. Файли формату csv.%enditem
\par
%beginitem
2. Написати функцію, що в текстовому файлі замінює всі входження фрази \textbf{не зараховано} на слово UPS. Кількість білих символів між частинами фрази може бути довільна (і там можуть бути не тільки пробіли). Результат має бути записаний у новий файл.

%enditem
\par
%beginitem
3. Текстовий файл містить послідовність цілих чисел $a_1$, $a_2$, $\ldots$, $a_t$, записану у форматі одне число на рядок. Знайти тринадцяте найбільше значення, що записане у файлі. Кожне значення, що записане у файлі, враховується тільки один раз. Кількість чисел у файлі така, що вони не можуть бути завантажені повністю в пам'ять програми. Розв'язання оформити у вигляді функції.

%enditem
}
{\smallskip\smallskip \begin {scriptsize} Затверджено на засіданні кафедри теоретичної кібернетики, протокол \No 4  від   25.01.2024 р.\par\smallskip\smallskip\smallskip\smallskip
{\parbox {6.5 cm} {Зав. кафедри \dotfill Ю.В.~Крак}}\hspace {0.7cm} {\hbox to 6.5 cm { Екзаменатор \dotfill Т.О.~Карнаух}} \end{scriptsize}}
%end
\vspace{0.9cm}
%begin
{{\begin {scriptsize}\par
{ \center  Київський національний університет імені Тараса Шевченка \par}
{\parbox {5 cm}{Освітня програма \hfill}{\parbox {7 cm}{Прикладна математика\hfill}}\parbox {6 cm}{\hfill Семестр 2}}\par
{\parbox {5 cm} {Навчальний предмет \hfill}\parbox {7 cm}{Програмування \hfill}\parbox {6cm}{\hfill}\par}\vspace{-15pt} \end{scriptsize}}{\center Екзаменаційний білет \No \textbf { 9 }\par}}
{%beginitem
1. Множини.%enditem
\par
%beginitem
2. Написати функцію, що за текстовим файлом, що містить зображення цілих чисел, дійсних чисел та точок площини, розділені білими символами, обчислює кількість входжень у файл зображень точок. Зображення точки має вигляд: відкриваюча дужка, x-координата, кома, y-координата, закриваюча дужка, наприклад, ( 13.5 , 45.1E-1) або ( 14, -3) . Білі символи можуть зустрічатись у довільній кількості. Файл дуже великий і його вміст завантажений у пам'ять бути не може.

%enditem
\par
%beginitem
3. Текстовий файл містить послідовність цілих чисел $a_1$, $a_2$, $\ldots$, $a_t$, записану у форматі одне число на рядок. Знайти тринадцяте найменше значення, що записане у файлі. Кожне значення, що записане у файлі, враховується тільки один раз. Кількість чисел у файлі така, що вони не можуть бути завантажені повністю в пам'ять програми. Розв'язання оформити у вигляді функції.

%enditem
}
{\smallskip\smallskip \begin {scriptsize} Затверджено на засіданні кафедри теоретичної кібернетики, протокол \No 4  від   25.01.2024 р.\par\smallskip\smallskip\smallskip\smallskip
{\parbox {6.5 cm} {Зав. кафедри \dotfill Ю.В.~Крак}}\hspace {0.7cm} {\hbox to 6.5 cm { Екзаменатор \dotfill Т.О.~Карнаух}} \end{scriptsize}}
%end
\newpage
%begin
{{\begin {scriptsize}\par
{ \center  Київський національний університет імені Тараса Шевченка \par}
{\parbox {5 cm}{Освітня програма \hfill}{\parbox {7 cm}{Прикладна математика\hfill}}\parbox {6 cm}{\hfill Семестр 2}}\par
{\parbox {5 cm} {Навчальний предмет \hfill}\parbox {7 cm}{Програмування \hfill}\parbox {6cm}{\hfill}\par}\vspace{-15pt} \end{scriptsize}}{\center Екзаменаційний білет \No \textbf { 10 }\par}}
{%beginitem
1. Емпіричний підхід до визначення складності обчислень.%enditem
\par
%beginitem
2. Написати функцію, що за текстовим файлом з статтею визначає скільки разів у ньому містяться слова, що починаються на Київ та Kyiv.  Переноси слів між рядками у файлі відсутні. Файл дуже великий і його вміст завантажений у пам'ять бути не може.

%enditem
\par
%beginitem
3. Текстовий файл містить послідовність цілих чисел $a_1$, $a_2$, $\ldots$, $a_t$, записану у форматі одне число на рядок. Знайти тринадцяте найменше значення, що записане у файлі. Кожне значення, що записане у файлі, враховується тільки один раз. Кількість чисел у файлі така, що вони не можуть бути завантажені повністю в пам'ять програми. Розв'язання оформити у вигляді функції.

%enditem
}
{\smallskip\smallskip \begin {scriptsize} Затверджено на засіданні кафедри теоретичної кібернетики, протокол \No 4  від   25.01.2024 р.\par\smallskip\smallskip\smallskip\smallskip
{\parbox {6.5 cm} {Зав. кафедри \dotfill Ю.В.~Крак}}\hspace {0.7cm} {\hbox to 6.5 cm { Екзаменатор \dotfill Т.О.~Карнаух}} \end{scriptsize}}
%end
\vspace{0.9cm}
%begin
{{\begin {scriptsize}\par
{ \center  Київський національний університет імені Тараса Шевченка \par}
{\parbox {5 cm}{Освітня програма \hfill}{\parbox {7 cm}{Прикладна математика\hfill}}\parbox {6 cm}{\hfill Семестр 2}}\par
{\parbox {5 cm} {Навчальний предмет \hfill}\parbox {7 cm}{Програмування \hfill}\parbox {6cm}{\hfill}\par}\vspace{-15pt} \end{scriptsize}}{\center Екзаменаційний білет \No \textbf { 11 }\par}}
{%beginitem
1. Ітератори.%enditem
\par
%beginitem
2. Написати функцію, що в текстовому файлі замінює всі зазначення ваги у вигляді \textbf{якесь ціле число кг} на слово WEIGHT (приклади зазначення ваги: 13 кг, 26   кг). Файл дуже великий і його вміст завантажений у пам'ять бути не може. Результат має бути записаний у новий файл.

%enditem
\par
%beginitem
3. Реалізувати операцію різниці множин з повтореннями. Множина з повтореннями задається переліком своїх елементів, записаних за неспаданням у текстовому файлі через білі символи. Результат у такому ж самому форматі (і теж за неспаданням) має бути записаний у текстовий файл. Питома функція має отримувати три шляхи до файлів. Розмір вхідних множин такий, що вони не можуть бути завантажені повністю в пам’ять програми.

%enditem
}
{\smallskip\smallskip \begin {scriptsize} Затверджено на засіданні кафедри теоретичної кібернетики, протокол \No 4  від   25.01.2024 р.\par\smallskip\smallskip\smallskip\smallskip
{\parbox {6.5 cm} {Зав. кафедри \dotfill Ю.В.~Крак}}\hspace {0.7cm} {\hbox to 6.5 cm { Екзаменатор \dotfill Т.О.~Карнаух}} \end{scriptsize}}
%end
\vspace{0.9cm}
%begin
{{\begin {scriptsize}\par
{ \center  Київський національний університет імені Тараса Шевченка \par}
{\parbox {5 cm}{Освітня програма \hfill}{\parbox {7 cm}{Прикладна математика\hfill}}\parbox {6 cm}{\hfill Семестр 2}}\par
{\parbox {5 cm} {Навчальний предмет \hfill}\parbox {7 cm}{Програмування \hfill}\parbox {6cm}{\hfill}\par}\vspace{-15pt} \end{scriptsize}}{\center Екзаменаційний білет \No \textbf { 12 }\par}}
{%beginitem
1. Лінійні зв’язані структури даних.%enditem
\par
%beginitem
2. Написати функцію, що для текстового файлу, що містить послідовність точок площини, обчислює центр мас заданих точок (центр мас -- точка, координати якої є середнім арифметичним відповідних координат). Для кожної точки на окремому рядку наводиться її зображення: відкриваюча дужка, x-координата, кома, y-координата, закриваюча дужка, наприклад, ( 13.5 , 45.1E-1). Білі символи можуть зустрічатись у довільній кількості. Файл дуже великий і його вміст завантажений у пам'ять бути не може.

%enditem
\par
%beginitem
3. Текстовий файл містить послідовність цілих чисел $a_1$, $a_2$, $\ldots$, $a_t$, записану у форматі одне число на рядок. Знайти сімнадцяте найбільше значення, що записане у файлі. Кожне значення, що записане у файлі, враховується стільки разів, скільки входить у файл. Кількість чисел у файлі така, що вони не можуть бути завантажені повністю в пам'ять програми. Розв'язання оформити у вигляді функції.

%enditem
}
{\smallskip\smallskip \begin {scriptsize} Затверджено на засіданні кафедри теоретичної кібернетики, протокол \No 4  від   25.01.2024 р.\par\smallskip\smallskip\smallskip\smallskip
{\parbox {6.5 cm} {Зав. кафедри \dotfill Ю.В.~Крак}}\hspace {0.7cm} {\hbox to 6.5 cm { Екзаменатор \dotfill Т.О.~Карнаух}} \end{scriptsize}}
%end
\newpage
%begin
{{\begin {scriptsize}\par
{ \center  Київський національний університет імені Тараса Шевченка \par}
{\parbox {5 cm}{Освітня програма \hfill}{\parbox {7 cm}{Прикладна математика\hfill}}\parbox {6 cm}{\hfill Семестр 2}}\par
{\parbox {5 cm} {Навчальний предмет \hfill}\parbox {7 cm}{Програмування \hfill}\parbox {6cm}{\hfill}\par}\vspace{-15pt} \end{scriptsize}}{\center Екзаменаційний білет \No \textbf { 13 }\par}}
{%beginitem
1. Словники.%enditem
\par
%beginitem
2. Написати функцію, що за текстовим файлом з статтею визначає  скільки разів у ньому містяться слова, що закінчуються на \textbf{ський} або на \textbf{зький} (рід, число та відмінок значення не мають, тобто \textit{міська}, \textit{міському} тощо теж мають рахуватись). Файл дуже великий і його вміст завантажений у пам'ять бути не може.

%enditem
\par
%beginitem
3. Текстовий файл містить послідовність цілих чисел $a_1$, $a_2$, $\ldots$, $a_t$, записану у форматі одне число на рядок. Знайти сімнадцяте найменше значення, що записане у файлі. Кожне значення, що записане у файлі, враховується стільки разів, скільки входить у файл. Кількість чисел у файлі така, що вони не можуть бути завантажені повністю в пам'ять програми. Розв'язання оформити у вигляді функції.

%enditem
}
{\smallskip\smallskip \begin {scriptsize} Затверджено на засіданні кафедри теоретичної кібернетики, протокол \No 4  від   25.01.2024 р.\par\smallskip\smallskip\smallskip\smallskip
{\parbox {6.5 cm} {Зав. кафедри \dotfill Ю.В.~Крак}}\hspace {0.7cm} {\hbox to 6.5 cm { Екзаменатор \dotfill Т.О.~Карнаух}} \end{scriptsize}}
%end
\vspace{0.9cm}
%begin
{{\begin {scriptsize}\par
{ \center  Київський національний університет імені Тараса Шевченка \par}
{\parbox {5 cm}{Освітня програма \hfill}{\parbox {7 cm}{Прикладна математика\hfill}}\parbox {6 cm}{\hfill Семестр 2}}\par
{\parbox {5 cm} {Навчальний предмет \hfill}\parbox {7 cm}{Програмування \hfill}\parbox {6cm}{\hfill}\par}\vspace{-15pt} \end{scriptsize}}{\center Екзаменаційний білет \No \textbf { 14 }\par}}
{%beginitem
1. Статичні поля, статичні методи, методи класів.%enditem
\par
%beginitem
2. Написати функцію, що в текстовому файлі замінює всі цілі числа на слово NUMBER. Файл дуже великий і його вміст завантажений у пам'ять бути не може. Результат має бути записаний у новий файл.

%enditem
\par
%beginitem
3. Реалізувати операцію перетину множин з повтореннями. Множина з повтореннями задається переліком своїх елементів, записаних за неспаданням у текстовому файлі через білі символи. Результат у такому ж самому форматі (і теж за неспаданням) має бути записаний у текстовий файл. Питома функція має отримувати три шляхи до файлів. Розмір вхідних множин такий, що вони не можуть бути завантажені повністю в пам’ять програми.

%enditem
}
{\smallskip\smallskip \begin {scriptsize} Затверджено на засіданні кафедри теоретичної кібернетики, протокол \No 4  від   25.01.2024 р.\par\smallskip\smallskip\smallskip\smallskip
{\parbox {6.5 cm} {Зав. кафедри \dotfill Ю.В.~Крак}}\hspace {0.7cm} {\hbox to 6.5 cm { Екзаменатор \dotfill Т.О.~Карнаух}} \end{scriptsize}}
%end
\vspace{0.9cm}
%begin
{{\begin {scriptsize}\par
{ \center  Київський національний університет імені Тараса Шевченка \par}
{\parbox {5 cm}{Освітня програма \hfill}{\parbox {7 cm}{Прикладна математика\hfill}}\parbox {6 cm}{\hfill Семестр 2}}\par
{\parbox {5 cm} {Навчальний предмет \hfill}\parbox {7 cm}{Програмування \hfill}\parbox {6cm}{\hfill}\par}\vspace{-15pt} \end{scriptsize}}{\center Екзаменаційний білет \No \textbf { 15 }\par}}
{%beginitem
1. Дерева. Обходи дерев.%enditem
\par
%beginitem
2. Написати функцію, що в текстовому файлі замінює всі входження фрази \textbf{не зараховано} на слово UPS. Кількість білих символів між частинами фрази може бути довільна (і там можуть бути не тільки пробіли). Результат має бути записаний у новий файл.

%enditem
\par
%beginitem
3. Кожен з двох текстових файлів містить послідовність дійсних чисел, записану у форматі одне число на рядок.  Перевірити, чи задають вони одну ту саму послідовність. Кількість чисел у кожному з двох файлів така, що вони не можуть бути завантажені повністю в пам'ять програми. Розв'язання оформити у вигляді функції.

%enditem
}
{\smallskip\smallskip \begin {scriptsize} Затверджено на засіданні кафедри теоретичної кібернетики, протокол \No 4  від   25.01.2024 р.\par\smallskip\smallskip\smallskip\smallskip
{\parbox {6.5 cm} {Зав. кафедри \dotfill Ю.В.~Крак}}\hspace {0.7cm} {\hbox to 6.5 cm { Екзаменатор \dotfill Т.О.~Карнаух}} \end{scriptsize}}
%end
\newpage
%begin
{{\begin {scriptsize}\par
{ \center  Київський національний університет імені Тараса Шевченка \par}
{\parbox {5 cm}{Освітня програма \hfill}{\parbox {7 cm}{Прикладна математика\hfill}}\parbox {6 cm}{\hfill Семестр 2}}\par
{\parbox {5 cm} {Навчальний предмет \hfill}\parbox {7 cm}{Програмування \hfill}\parbox {6cm}{\hfill}\par}\vspace{-15pt} \end{scriptsize}}{\center Екзаменаційний білет \No \textbf { 16 }\par}}
{%beginitem
1. Обробка файлів: режими відкриття файлів.%enditem
\par
%beginitem
2. Написати функцію, що для текстового файлу, що містить послідовність точок площини, обчислює центр мас заданих точок (центр мас -- точка, координати якої є середнім арифметичним відповідних координат). Для кожної точки на окремому рядку наводиться її зображення: відкриваюча дужка, x-координата, кома, y-координата, закриваюча дужка, наприклад, ( 13.5 , 45.1E-1). Білі символи можуть зустрічатись у довільній кількості. Файл дуже великий і його вміст завантажений у пам'ять бути не може.

%enditem
\par
%beginitem
3. Реалізувати операцію симетричної різниці множин з повтореннями. Множина з повтореннями задається переліком своїх елементів, записаних за неспаданням у текстовому файлі через білі символи. Результат у такому ж самому форматі (і теж за неспаданням) має бути записаний у текстовий файл. Питома функція має отримувати три шляхи до файлів. Розмір вхідних множин такий, що вони не можуть бути завантажені повністю в пам’ять програми.

%enditem
}
{\smallskip\smallskip \begin {scriptsize} Затверджено на засіданні кафедри теоретичної кібернетики, протокол \No 4  від   25.01.2024 р.\par\smallskip\smallskip\smallskip\smallskip
{\parbox {6.5 cm} {Зав. кафедри \dotfill Ю.В.~Крак}}\hspace {0.7cm} {\hbox to 6.5 cm { Екзаменатор \dotfill Т.О.~Карнаух}} \end{scriptsize}}
%end
\vspace{0.9cm}
%begin
{{\begin {scriptsize}\par
{ \center  Київський національний університет імені Тараса Шевченка \par}
{\parbox {5 cm}{Освітня програма \hfill}{\parbox {7 cm}{Прикладна математика\hfill}}\parbox {6 cm}{\hfill Семестр 2}}\par
{\parbox {5 cm} {Навчальний предмет \hfill}\parbox {7 cm}{Програмування \hfill}\parbox {6cm}{\hfill}\par}\vspace{-15pt} \end{scriptsize}}{\center Екзаменаційний білет \No \textbf { 17 }\par}}
{%beginitem
1. Математичний підхід до визначення складності обчислень.%enditem
\par
%beginitem
2. Написати функцію, що в текстовому файлі замінює всі цілі числа на слово NUMBER. Файл дуже великий і його вміст завантажений у пам'ять бути не може. Результат має бути записаний у новий файл.

%enditem
\par
%beginitem
3. Реалізувати операцію злиття множин з повтореннями. Множина з повтореннями задається переліком своїх елементів, записаних за неспаданням у текстовому файлі через білі символи. Результат у такому ж самому форматі (і теж за неспаданням) має бути записаний у текстовий файл. Питома функція має отримувати три шляхи до файлів. Розмір вхідних множин такий, що вони не можуть бути завантажені повністю в пам’ять програми.

%enditem
}
{\smallskip\smallskip \begin {scriptsize} Затверджено на засіданні кафедри теоретичної кібернетики, протокол \No 4  від   25.01.2024 р.\par\smallskip\smallskip\smallskip\smallskip
{\parbox {6.5 cm} {Зав. кафедри \dotfill Ю.В.~Крак}}\hspace {0.7cm} {\hbox to 6.5 cm { Екзаменатор \dotfill Т.О.~Карнаух}} \end{scriptsize}}
%end
\vspace{0.9cm}
%begin
{{\begin {scriptsize}\par
{ \center  Київський національний університет імені Тараса Шевченка \par}
{\parbox {5 cm}{Освітня програма \hfill}{\parbox {7 cm}{Прикладна математика\hfill}}\parbox {6 cm}{\hfill Семестр 2}}\par
{\parbox {5 cm} {Навчальний предмет \hfill}\parbox {7 cm}{Програмування \hfill}\parbox {6cm}{\hfill}\par}\vspace{-15pt} \end{scriptsize}}{\center Екзаменаційний білет \No \textbf { 18 }\par}}
{%beginitem
1. Поняття часової та просторової складності обчислень. Складність у середньому. Асимптотичні оцінки складності.%enditem
\par
%beginitem
2. Написати функцію, що за текстовим файлом, що містить зображення цілих чисел, дійсних чисел та точок площини, розділені білими символами, обчислює кількість входжень у файл зображень точок. Зображення точки має вигляд: відкриваюча дужка, x-координата, кома, y-координата, закриваюча дужка, наприклад, ( 13.5 , 45.1E-1) або ( 14, -3) . Білі символи можуть зустрічатись у довільній кількості. Файл дуже великий і його вміст завантажений у пам'ять бути не може.

%enditem
\par
%beginitem
3. Текстовий файл містить послідовність цілих чисел $a_1$, $a_2$, $\ldots$, $a_t$, записану у форматі одне число на рядок. Знайти тринадцяте найбільше значення, що записане у файлі. Кожне значення, що записане у файлі, враховується тільки один раз. Кількість чисел у файлі така, що вони не можуть бути завантажені повністю в пам'ять програми. Розв'язання оформити у вигляді функції.

%enditem
}
{\smallskip\smallskip \begin {scriptsize} Затверджено на засіданні кафедри теоретичної кібернетики, протокол \No 4  від   25.01.2024 р.\par\smallskip\smallskip\smallskip\smallskip
{\parbox {6.5 cm} {Зав. кафедри \dotfill Ю.В.~Крак}}\hspace {0.7cm} {\hbox to 6.5 cm { Екзаменатор \dotfill Т.О.~Карнаух}} \end{scriptsize}}
%end
\newpage
%begin
{{\begin {scriptsize}\par
{ \center  Київський національний університет імені Тараса Шевченка \par}
{\parbox {5 cm}{Освітня програма \hfill}{\parbox {7 cm}{Прикладна математика\hfill}}\parbox {6 cm}{\hfill Семестр 2}}\par
{\parbox {5 cm} {Навчальний предмет \hfill}\parbox {7 cm}{Програмування \hfill}\parbox {6cm}{\hfill}\par}\vspace{-15pt} \end{scriptsize}}{\center Екзаменаційний білет \No \textbf { 19 }\par}}
{%beginitem
1. Емпіричний підхід до визначення складності обчислень.%enditem
\par
%beginitem
2. Написати функцію, що за текстовим файлом з статтею визначає скільки разів у ньому містяться слова, що починаються на Київ та Kyiv.  Переноси слів між рядками у файлі відсутні. Файл дуже великий і його вміст завантажений у пам'ять бути не може.

%enditem
\par
%beginitem
3. Кожен з двох текстових файлів містить послідовність дійсних чисел, записану у форматі одне число на рядок.  Перевірити, чи задають вони одну ту саму послідовність. Кількість чисел у кожному з двох файлів така, що вони не можуть бути завантажені повністю в пам'ять програми. Розв'язання оформити у вигляді функції.

%enditem
}
{\smallskip\smallskip \begin {scriptsize} Затверджено на засіданні кафедри теоретичної кібернетики, протокол \No 4  від   25.01.2024 р.\par\smallskip\smallskip\smallskip\smallskip
{\parbox {6.5 cm} {Зав. кафедри \dotfill Ю.В.~Крак}}\hspace {0.7cm} {\hbox to 6.5 cm { Екзаменатор \dotfill Т.О.~Карнаух}} \end{scriptsize}}
%end
\vspace{0.9cm}
%begin
{{\begin {scriptsize}\par
{ \center  Київський національний університет імені Тараса Шевченка \par}
{\parbox {5 cm}{Освітня програма \hfill}{\parbox {7 cm}{Прикладна математика\hfill}}\parbox {6 cm}{\hfill Семестр 2}}\par
{\parbox {5 cm} {Навчальний предмет \hfill}\parbox {7 cm}{Програмування \hfill}\parbox {6cm}{\hfill}\par}\vspace{-15pt} \end{scriptsize}}{\center Екзаменаційний білет \No \textbf { 20 }\par}}
{%beginitem
1. Словники.%enditem
\par
%beginitem
2. Написати функцію, що за текстовим файлом з статтею визначає  скільки разів у ньому містяться слова, що закінчуються на \textbf{ський} або на \textbf{зький} (рід, число та відмінок значення не мають, тобто \textit{міська}, \textit{міському} тощо теж мають рахуватись). Файл дуже великий і його вміст завантажений у пам'ять бути не може.

%enditem
\par
%beginitem
3. Реалізувати операцію злиття множин з повтореннями. Множина з повтореннями задається переліком своїх елементів, записаних за неспаданням у текстовому файлі через білі символи. Результат у такому ж самому форматі (і теж за неспаданням) має бути записаний у текстовий файл. Питома функція має отримувати три шляхи до файлів. Розмір вхідних множин такий, що вони не можуть бути завантажені повністю в пам’ять програми.

%enditem
}
{\smallskip\smallskip \begin {scriptsize} Затверджено на засіданні кафедри теоретичної кібернетики, протокол \No 4  від   25.01.2024 р.\par\smallskip\smallskip\smallskip\smallskip
{\parbox {6.5 cm} {Зав. кафедри \dotfill Ю.В.~Крак}}\hspace {0.7cm} {\hbox to 6.5 cm { Екзаменатор \dotfill Т.О.~Карнаух}} \end{scriptsize}}
%end
\vspace{0.9cm}
%begin
{{\begin {scriptsize}\par
{ \center  Київський національний університет імені Тараса Шевченка \par}
{\parbox {5 cm}{Освітня програма \hfill}{\parbox {7 cm}{Прикладна математика\hfill}}\parbox {6 cm}{\hfill Семестр 2}}\par
{\parbox {5 cm} {Навчальний предмет \hfill}\parbox {7 cm}{Програмування \hfill}\parbox {6cm}{\hfill}\par}\vspace{-15pt} \end{scriptsize}}{\center Екзаменаційний білет \No \textbf { 21 }\par}}
{%beginitem
1. Алгоритм пошуку в ширину та задачі, які він може вирішувати.%enditem
\par
%beginitem
2. Написати функцію, що в текстовому файлі замінює всі зазначення ваги у вигляді \textbf{якесь ціле число кг} на слово WEIGHT (приклади зазначення ваги: 13 кг, 26   кг). Файл дуже великий і його вміст завантажений у пам'ять бути не може. Результат має бути записаний у новий файл.

%enditem
\par
%beginitem
3. Текстовий файл містить послідовність цілих чисел $a_1$, $a_2$, $\ldots$, $a_t$, записану у форматі одне число на рядок. Знайти сімнадцяте найбільше значення, що записане у файлі. Кожне значення, що записане у файлі, враховується стільки разів, скільки входить у файл. Кількість чисел у файлі така, що вони не можуть бути завантажені повністю в пам'ять програми. Розв'язання оформити у вигляді функції.

%enditem
}
{\smallskip\smallskip \begin {scriptsize} Затверджено на засіданні кафедри теоретичної кібернетики, протокол \No 4  від   25.01.2024 р.\par\smallskip\smallskip\smallskip\smallskip
{\parbox {6.5 cm} {Зав. кафедри \dotfill Ю.В.~Крак}}\hspace {0.7cm} {\hbox to 6.5 cm { Екзаменатор \dotfill Т.О.~Карнаух}} \end{scriptsize}}
%end
\newpage
%begin
{{\begin {scriptsize}\par
{ \center  Київський національний університет імені Тараса Шевченка \par}
{\parbox {5 cm}{Освітня програма \hfill}{\parbox {7 cm}{Прикладна математика\hfill}}\parbox {6 cm}{\hfill Семестр 2}}\par
{\parbox {5 cm} {Навчальний предмет \hfill}\parbox {7 cm}{Програмування \hfill}\parbox {6cm}{\hfill}\par}\vspace{-15pt} \end{scriptsize}}{\center Екзаменаційний білет \No \textbf { 22 }\par}}
{%beginitem
1. Математичний підхід до визначення складності обчислень.%enditem
\par
%beginitem
2. Написати функцію, що для текстового файлу, що містить послідовність точок площини, обчислює центр мас заданих точок (центр мас -- точка, координати якої є середнім арифметичним відповідних координат). Для кожної точки на окремому рядку наводиться її зображення: відкриваюча дужка, x-координата, кома, y-координата, закриваюча дужка, наприклад, ( 13.5 , 45.1E-1). Білі символи можуть зустрічатись у довільній кількості. Файл дуже великий і його вміст завантажений у пам'ять бути не може.

%enditem
\par
%beginitem
3. Текстовий файл містить послідовність цілих чисел $a_1$, $a_2$, $\ldots$, $a_t$, записану у форматі одне число на рядок. Знайти тринадцяте найменше значення, що записане у файлі. Кожне значення, що записане у файлі, враховується тільки один раз. Кількість чисел у файлі така, що вони не можуть бути завантажені повністю в пам'ять програми. Розв'язання оформити у вигляді функції.

%enditem
}
{\smallskip\smallskip \begin {scriptsize} Затверджено на засіданні кафедри теоретичної кібернетики, протокол \No 4  від   25.01.2024 р.\par\smallskip\smallskip\smallskip\smallskip
{\parbox {6.5 cm} {Зав. кафедри \dotfill Ю.В.~Крак}}\hspace {0.7cm} {\hbox to 6.5 cm { Екзаменатор \dotfill Т.О.~Карнаух}} \end{scriptsize}}
%end
\vspace{0.9cm}
%begin
{{\begin {scriptsize}\par
{ \center  Київський національний університет імені Тараса Шевченка \par}
{\parbox {5 cm}{Освітня програма \hfill}{\parbox {7 cm}{Прикладна математика\hfill}}\parbox {6 cm}{\hfill Семестр 2}}\par
{\parbox {5 cm} {Навчальний предмет \hfill}\parbox {7 cm}{Програмування \hfill}\parbox {6cm}{\hfill}\par}\vspace{-15pt} \end{scriptsize}}{\center Екзаменаційний білет \No \textbf { 23 }\par}}
{%beginitem
1. Множини.%enditem
\par
%beginitem
2. Написати функцію, що за текстовим файлом з статтею визначає  скільки разів у ньому містяться слова, що закінчуються на \textbf{ський} або на \textbf{зький} (рід, число та відмінок значення не мають, тобто \textit{міська}, \textit{міському} тощо теж мають рахуватись). Файл дуже великий і його вміст завантажений у пам'ять бути не може.

%enditem
\par
%beginitem
3. Реалізувати операцію симетричної різниці множин з повтореннями. Множина з повтореннями задається переліком своїх елементів, записаних за неспаданням у текстовому файлі через білі символи. Результат у такому ж самому форматі (і теж за неспаданням) має бути записаний у текстовий файл. Питома функція має отримувати три шляхи до файлів. Розмір вхідних множин такий, що вони не можуть бути завантажені повністю в пам’ять програми.

%enditem
}
{\smallskip\smallskip \begin {scriptsize} Затверджено на засіданні кафедри теоретичної кібернетики, протокол \No 4  від   25.01.2024 р.\par\smallskip\smallskip\smallskip\smallskip
{\parbox {6.5 cm} {Зав. кафедри \dotfill Ю.В.~Крак}}\hspace {0.7cm} {\hbox to 6.5 cm { Екзаменатор \dotfill Т.О.~Карнаух}} \end{scriptsize}}
%end
\vspace{0.9cm}
%begin
{{\begin {scriptsize}\par
{ \center  Київський національний університет імені Тараса Шевченка \par}
{\parbox {5 cm}{Освітня програма \hfill}{\parbox {7 cm}{Прикладна математика\hfill}}\parbox {6 cm}{\hfill Семестр 2}}\par
{\parbox {5 cm} {Навчальний предмет \hfill}\parbox {7 cm}{Програмування \hfill}\parbox {6cm}{\hfill}\par}\vspace{-15pt} \end{scriptsize}}{\center Екзаменаційний білет \No \textbf { 24 }\par}}
{%beginitem
1. Багатовимірні масиви бібліотеки numpy: у чому відмінність моделювання матриці багатовимірним масивом бібліотеки numpy та як список списків.%enditem
\par
%beginitem
2. Написати функцію, що в текстовому файлі замінює всі зазначення ваги у вигляді \textbf{якесь ціле число кг} на слово WEIGHT (приклади зазначення ваги: 13 кг, 26   кг). Файл дуже великий і його вміст завантажений у пам'ять бути не може. Результат має бути записаний у новий файл.

%enditem
\par
%beginitem
3. Текстовий файл містить послідовність цілих чисел $a_1$, $a_2$, $\ldots$, $a_t$, записану у форматі одне число на рядок. Знайти сімнадцяте найменше значення, що записане у файлі. Кожне значення, що записане у файлі, враховується стільки разів, скільки входить у файл. Кількість чисел у файлі така, що вони не можуть бути завантажені повністю в пам'ять програми. Розв'язання оформити у вигляді функції.

%enditem
}
{\smallskip\smallskip \begin {scriptsize} Затверджено на засіданні кафедри теоретичної кібернетики, протокол \No 4  від   25.01.2024 р.\par\smallskip\smallskip\smallskip\smallskip
{\parbox {6.5 cm} {Зав. кафедри \dotfill Ю.В.~Крак}}\hspace {0.7cm} {\hbox to 6.5 cm { Екзаменатор \dotfill Т.О.~Карнаух}} \end{scriptsize}}
%end
\newpage
%begin
{{\begin {scriptsize}\par
{ \center  Київський національний університет імені Тараса Шевченка \par}
{\parbox {5 cm}{Освітня програма \hfill}{\parbox {7 cm}{Прикладна математика\hfill}}\parbox {6 cm}{\hfill Семестр 2}}\par
{\parbox {5 cm} {Навчальний предмет \hfill}\parbox {7 cm}{Програмування \hfill}\parbox {6cm}{\hfill}\par}\vspace{-15pt} \end{scriptsize}}{\center Екзаменаційний білет \No \textbf { 25 }\par}}
{%beginitem
1. Ітератори.%enditem
\par
%beginitem
2. Написати функцію, що в текстовому файлі замінює всі входження фрази \textbf{не зараховано} на слово UPS. Кількість білих символів між частинами фрази може бути довільна (і там можуть бути не тільки пробіли). Результат має бути записаний у новий файл.

%enditem
\par
%beginitem
3. Реалізувати операцію різниці множин з повтореннями. Множина з повтореннями задається переліком своїх елементів, записаних за неспаданням у текстовому файлі через білі символи. Результат у такому ж самому форматі (і теж за неспаданням) має бути записаний у текстовий файл. Питома функція має отримувати три шляхи до файлів. Розмір вхідних множин такий, що вони не можуть бути завантажені повністю в пам’ять програми.

%enditem
}
{\smallskip\smallskip \begin {scriptsize} Затверджено на засіданні кафедри теоретичної кібернетики, протокол \No 4  від   25.01.2024 р.\par\smallskip\smallskip\smallskip\smallskip
{\parbox {6.5 cm} {Зав. кафедри \dotfill Ю.В.~Крак}}\hspace {0.7cm} {\hbox to 6.5 cm { Екзаменатор \dotfill Т.О.~Карнаух}} \end{scriptsize}}
%end
\vspace{0.9cm}
%begin
{{\begin {scriptsize}\par
{ \center  Київський національний університет імені Тараса Шевченка \par}
{\parbox {5 cm}{Освітня програма \hfill}{\parbox {7 cm}{Прикладна математика\hfill}}\parbox {6 cm}{\hfill Семестр 2}}\par
{\parbox {5 cm} {Навчальний предмет \hfill}\parbox {7 cm}{Програмування \hfill}\parbox {6cm}{\hfill}\par}\vspace{-15pt} \end{scriptsize}}{\center Екзаменаційний білет \No \textbf { 26 }\par}}
{%beginitem
1. Файли формату csv.%enditem
\par
%beginitem
2. Написати функцію, що за текстовим файлом з статтею визначає скільки разів у ньому містяться слова, що починаються на Київ та Kyiv.  Переноси слів між рядками у файлі відсутні. Файл дуже великий і його вміст завантажений у пам'ять бути не може.

%enditem
\par
%beginitem
3. Реалізувати операцію перетину множин з повтореннями. Множина з повтореннями задається переліком своїх елементів, записаних за неспаданням у текстовому файлі через білі символи. Результат у такому ж самому форматі (і теж за неспаданням) має бути записаний у текстовий файл. Питома функція має отримувати три шляхи до файлів. Розмір вхідних множин такий, що вони не можуть бути завантажені повністю в пам’ять програми.

%enditem
}
{\smallskip\smallskip \begin {scriptsize} Затверджено на засіданні кафедри теоретичної кібернетики, протокол \No 4  від   25.01.2024 р.\par\smallskip\smallskip\smallskip\smallskip
{\parbox {6.5 cm} {Зав. кафедри \dotfill Ю.В.~Крак}}\hspace {0.7cm} {\hbox to 6.5 cm { Екзаменатор \dotfill Т.О.~Карнаух}} \end{scriptsize}}
%end
\vspace{0.9cm}
%begin
{{\begin {scriptsize}\par
{ \center  Київський національний університет імені Тараса Шевченка \par}
{\parbox {5 cm}{Освітня програма \hfill}{\parbox {7 cm}{Прикладна математика\hfill}}\parbox {6 cm}{\hfill Семестр 2}}\par
{\parbox {5 cm} {Навчальний предмет \hfill}\parbox {7 cm}{Програмування \hfill}\parbox {6cm}{\hfill}\par}\vspace{-15pt} \end{scriptsize}}{\center Екзаменаційний білет \No \textbf { 27 }\par}}
{%beginitem
1. Файли формату json.%enditem
\par
%beginitem
2. Написати функцію, що за текстовим файлом, що містить зображення цілих чисел, дійсних чисел та точок площини, розділені білими символами, обчислює кількість входжень у файл зображень точок. Зображення точки має вигляд: відкриваюча дужка, x-координата, кома, y-координата, закриваюча дужка, наприклад, ( 13.5 , 45.1E-1) або ( 14, -3) . Білі символи можуть зустрічатись у довільній кількості. Файл дуже великий і його вміст завантажений у пам'ять бути не може.

%enditem
\par
%beginitem
3. Текстовий файл містить послідовність цілих чисел $a_1$, $a_2$, $\ldots$, $a_t$, записану у форматі одне число на рядок. Знайти тринадцяте найбільше значення, що записане у файлі. Кожне значення, що записане у файлі, враховується тільки один раз. Кількість чисел у файлі така, що вони не можуть бути завантажені повністю в пам'ять програми. Розв'язання оформити у вигляді функції.

%enditem
}
{\smallskip\smallskip \begin {scriptsize} Затверджено на засіданні кафедри теоретичної кібернетики, протокол \No 4  від   25.01.2024 р.\par\smallskip\smallskip\smallskip\smallskip
{\parbox {6.5 cm} {Зав. кафедри \dotfill Ю.В.~Крак}}\hspace {0.7cm} {\hbox to 6.5 cm { Екзаменатор \dotfill Т.О.~Карнаух}} \end{scriptsize}}
%end
\newpage
%begin
{{\begin {scriptsize}\par
{ \center  Київський національний університет імені Тараса Шевченка \par}
{\parbox {5 cm}{Освітня програма \hfill}{\parbox {7 cm}{Прикладна математика\hfill}}\parbox {6 cm}{\hfill Семестр 2}}\par
{\parbox {5 cm} {Навчальний предмет \hfill}\parbox {7 cm}{Програмування \hfill}\parbox {6cm}{\hfill}\par}\vspace{-15pt} \end{scriptsize}}{\center Екзаменаційний білет \No \textbf { 28 }\par}}
{%beginitem
1. Алгоритм пошуку в глибину та задачі, які він може вирішувати.%enditem
\par
%beginitem
2. Написати функцію, що в текстовому файлі замінює всі цілі числа на слово NUMBER. Файл дуже великий і його вміст завантажений у пам'ять бути не може. Результат має бути записаний у новий файл.

%enditem
\par
%beginitem
3. Реалізувати операцію злиття множин з повтореннями. Множина з повтореннями задається переліком своїх елементів, записаних за неспаданням у текстовому файлі через білі символи. Результат у такому ж самому форматі (і теж за неспаданням) має бути записаний у текстовий файл. Питома функція має отримувати три шляхи до файлів. Розмір вхідних множин такий, що вони не можуть бути завантажені повністю в пам’ять програми.

%enditem
}
{\smallskip\smallskip \begin {scriptsize} Затверджено на засіданні кафедри теоретичної кібернетики, протокол \No 4  від   25.01.2024 р.\par\smallskip\smallskip\smallskip\smallskip
{\parbox {6.5 cm} {Зав. кафедри \dotfill Ю.В.~Крак}}\hspace {0.7cm} {\hbox to 6.5 cm { Екзаменатор \dotfill Т.О.~Карнаух}} \end{scriptsize}}
%end
\vspace{0.9cm}
%begin
{{\begin {scriptsize}\par
{ \center  Київський національний університет імені Тараса Шевченка \par}
{\parbox {5 cm}{Освітня програма \hfill}{\parbox {7 cm}{Прикладна математика\hfill}}\parbox {6 cm}{\hfill Семестр 2}}\par
{\parbox {5 cm} {Навчальний предмет \hfill}\parbox {7 cm}{Програмування \hfill}\parbox {6cm}{\hfill}\par}\vspace{-15pt} \end{scriptsize}}{\center Екзаменаційний білет \No \textbf { 29 }\par}}
{%beginitem
1. Дерева. Обходи дерев.%enditem
\par
%beginitem
2. Написати функцію, що за текстовим файлом з статтею визначає скільки разів у ньому містяться слова, що починаються на Київ та Kyiv.  Переноси слів між рядками у файлі відсутні. Файл дуже великий і його вміст завантажений у пам'ять бути не може.

%enditem
\par
%beginitem
3. Реалізувати операцію різниці множин з повтореннями. Множина з повтореннями задається переліком своїх елементів, записаних за неспаданням у текстовому файлі через білі символи. Результат у такому ж самому форматі (і теж за неспаданням) має бути записаний у текстовий файл. Питома функція має отримувати три шляхи до файлів. Розмір вхідних множин такий, що вони не можуть бути завантажені повністю в пам’ять програми.

%enditem
}
{\smallskip\smallskip \begin {scriptsize} Затверджено на засіданні кафедри теоретичної кібернетики, протокол \No 4  від   25.01.2024 р.\par\smallskip\smallskip\smallskip\smallskip
{\parbox {6.5 cm} {Зав. кафедри \dotfill Ю.В.~Крак}}\hspace {0.7cm} {\hbox to 6.5 cm { Екзаменатор \dotfill Т.О.~Карнаух}} \end{scriptsize}}
%end
\vspace{0.9cm}
%begin
{{\begin {scriptsize}\par
{ \center  Київський національний університет імені Тараса Шевченка \par}
{\parbox {5 cm}{Освітня програма \hfill}{\parbox {7 cm}{Прикладна математика\hfill}}\parbox {6 cm}{\hfill Семестр 2}}\par
{\parbox {5 cm} {Навчальний предмет \hfill}\parbox {7 cm}{Програмування \hfill}\parbox {6cm}{\hfill}\par}\vspace{-15pt} \end{scriptsize}}{\center Екзаменаційний білет \No \textbf { 30 }\par}}
{%beginitem
1. Поняття стеку та черги. Основні операції зі стеками та чергами.%enditem
\par
%beginitem
2. Написати функцію, що в текстовому файлі замінює всі входження фрази \textbf{не зараховано} на слово UPS. Кількість білих символів між частинами фрази може бути довільна (і там можуть бути не тільки пробіли). Результат має бути записаний у новий файл.

%enditem
\par
%beginitem
3. Текстовий файл містить послідовність цілих чисел $a_1$, $a_2$, $\ldots$, $a_t$, записану у форматі одне число на рядок. Знайти тринадцяте найменше значення, що записане у файлі. Кожне значення, що записане у файлі, враховується тільки один раз. Кількість чисел у файлі така, що вони не можуть бути завантажені повністю в пам'ять програми. Розв'язання оформити у вигляді функції.

%enditem
}
{\smallskip\smallskip \begin {scriptsize} Затверджено на засіданні кафедри теоретичної кібернетики, протокол \No 4  від   25.01.2024 р.\par\smallskip\smallskip\smallskip\smallskip
{\parbox {6.5 cm} {Зав. кафедри \dotfill Ю.В.~Крак}}\hspace {0.7cm} {\hbox to 6.5 cm { Екзаменатор \dotfill Т.О.~Карнаух}} \end{scriptsize}}
%end
\newpage
%begin
{{\begin {scriptsize}\par
{ \center  Київський національний університет імені Тараса Шевченка \par}
{\parbox {5 cm}{Освітня програма \hfill}{\parbox {7 cm}{Прикладна математика\hfill}}\parbox {6 cm}{\hfill Семестр 2}}\par
{\parbox {5 cm} {Навчальний предмет \hfill}\parbox {7 cm}{Програмування \hfill}\parbox {6cm}{\hfill}\par}\vspace{-15pt} \end{scriptsize}}{\center Екзаменаційний білет \No \textbf { 31 }\par}}
{%beginitem
1. Статичні поля, статичні методи, методи класів.%enditem
\par
%beginitem
2. Написати функцію, що для текстового файлу, що містить послідовність точок площини, обчислює центр мас заданих точок (центр мас -- точка, координати якої є середнім арифметичним відповідних координат). Для кожної точки на окремому рядку наводиться її зображення: відкриваюча дужка, x-координата, кома, y-координата, закриваюча дужка, наприклад, ( 13.5 , 45.1E-1). Білі символи можуть зустрічатись у довільній кількості. Файл дуже великий і його вміст завантажений у пам'ять бути не може.

%enditem
\par
%beginitem
3. Реалізувати операцію симетричної різниці множин з повтореннями. Множина з повтореннями задається переліком своїх елементів, записаних за неспаданням у текстовому файлі через білі символи. Результат у такому ж самому форматі (і теж за неспаданням) має бути записаний у текстовий файл. Питома функція має отримувати три шляхи до файлів. Розмір вхідних множин такий, що вони не можуть бути завантажені повністю в пам’ять програми.

%enditem
}
{\smallskip\smallskip \begin {scriptsize} Затверджено на засіданні кафедри теоретичної кібернетики, протокол \No 4  від   25.01.2024 р.\par\smallskip\smallskip\smallskip\smallskip
{\parbox {6.5 cm} {Зав. кафедри \dotfill Ю.В.~Крак}}\hspace {0.7cm} {\hbox to 6.5 cm { Екзаменатор \dotfill Т.О.~Карнаух}} \end{scriptsize}}
%end
\vspace{0.9cm}
%begin
{{\begin {scriptsize}\par
{ \center  Київський національний університет імені Тараса Шевченка \par}
{\parbox {5 cm}{Освітня програма \hfill}{\parbox {7 cm}{Прикладна математика\hfill}}\parbox {6 cm}{\hfill Семестр 2}}\par
{\parbox {5 cm} {Навчальний предмет \hfill}\parbox {7 cm}{Програмування \hfill}\parbox {6cm}{\hfill}\par}\vspace{-15pt} \end{scriptsize}}{\center Екзаменаційний білет \No \textbf { 32 }\par}}
{%beginitem
1. Обробка файлів: режими відкриття файлів.%enditem
\par
%beginitem
2. Написати функцію, що за текстовим файлом, що містить зображення цілих чисел, дійсних чисел та точок площини, розділені білими символами, обчислює кількість входжень у файл зображень точок. Зображення точки має вигляд: відкриваюча дужка, x-координата, кома, y-координата, закриваюча дужка, наприклад, ( 13.5 , 45.1E-1) або ( 14, -3) . Білі символи можуть зустрічатись у довільній кількості. Файл дуже великий і його вміст завантажений у пам'ять бути не може.

%enditem
\par
%beginitem
3. Текстовий файл містить послідовність цілих чисел $a_1$, $a_2$, $\ldots$, $a_t$, записану у форматі одне число на рядок. Знайти тринадцяте найбільше значення, що записане у файлі. Кожне значення, що записане у файлі, враховується тільки один раз. Кількість чисел у файлі така, що вони не можуть бути завантажені повністю в пам'ять програми. Розв'язання оформити у вигляді функції.

%enditem
}
{\smallskip\smallskip \begin {scriptsize} Затверджено на засіданні кафедри теоретичної кібернетики, протокол \No 4  від   25.01.2024 р.\par\smallskip\smallskip\smallskip\smallskip
{\parbox {6.5 cm} {Зав. кафедри \dotfill Ю.В.~Крак}}\hspace {0.7cm} {\hbox to 6.5 cm { Екзаменатор \dotfill Т.О.~Карнаух}} \end{scriptsize}}
%end
\vspace{0.9cm}
%begin
{{\begin {scriptsize}\par
{ \center  Київський національний університет імені Тараса Шевченка \par}
{\parbox {5 cm}{Освітня програма \hfill}{\parbox {7 cm}{Прикладна математика\hfill}}\parbox {6 cm}{\hfill Семестр 2}}\par
{\parbox {5 cm} {Навчальний предмет \hfill}\parbox {7 cm}{Програмування \hfill}\parbox {6cm}{\hfill}\par}\vspace{-15pt} \end{scriptsize}}{\center Екзаменаційний білет \No \textbf { 33 }\par}}
{%beginitem
1. Контекстні менеджери.%enditem
\par
%beginitem
2. Написати функцію, що в текстовому файлі замінює всі зазначення ваги у вигляді \textbf{якесь ціле число кг} на слово WEIGHT (приклади зазначення ваги: 13 кг, 26   кг). Файл дуже великий і його вміст завантажений у пам'ять бути не може. Результат має бути записаний у новий файл.

%enditem
\par
%beginitem
3. Реалізувати операцію перетину множин з повтореннями. Множина з повтореннями задається переліком своїх елементів, записаних за неспаданням у текстовому файлі через білі символи. Результат у такому ж самому форматі (і теж за неспаданням) має бути записаний у текстовий файл. Питома функція має отримувати три шляхи до файлів. Розмір вхідних множин такий, що вони не можуть бути завантажені повністю в пам’ять програми.

%enditem
}
{\smallskip\smallskip \begin {scriptsize} Затверджено на засіданні кафедри теоретичної кібернетики, протокол \No 4  від   25.01.2024 р.\par\smallskip\smallskip\smallskip\smallskip
{\parbox {6.5 cm} {Зав. кафедри \dotfill Ю.В.~Крак}}\hspace {0.7cm} {\hbox to 6.5 cm { Екзаменатор \dotfill Т.О.~Карнаух}} \end{scriptsize}}
%end
\newpage
%begin
{{\begin {scriptsize}\par
{ \center  Київський національний університет імені Тараса Шевченка \par}
{\parbox {5 cm}{Освітня програма \hfill}{\parbox {7 cm}{Прикладна математика\hfill}}\parbox {6 cm}{\hfill Семестр 2}}\par
{\parbox {5 cm} {Навчальний предмет \hfill}\parbox {7 cm}{Програмування \hfill}\parbox {6cm}{\hfill}\par}\vspace{-15pt} \end{scriptsize}}{\center Екзаменаційний білет \No \textbf { 34 }\par}}
{%beginitem
1. Лінійні зв’язані структури даних.%enditem
\par
%beginitem
2. Написати функцію, що за текстовим файлом з статтею визначає  скільки разів у ньому містяться слова, що закінчуються на \textbf{ський} або на \textbf{зький} (рід, число та відмінок значення не мають, тобто \textit{міська}, \textit{міському} тощо теж мають рахуватись). Файл дуже великий і його вміст завантажений у пам'ять бути не може.

%enditem
\par
%beginitem
3. Текстовий файл містить послідовність цілих чисел $a_1$, $a_2$, $\ldots$, $a_t$, записану у форматі одне число на рядок. Знайти сімнадцяте найбільше значення, що записане у файлі. Кожне значення, що записане у файлі, враховується стільки разів, скільки входить у файл. Кількість чисел у файлі така, що вони не можуть бути завантажені повністю в пам'ять програми. Розв'язання оформити у вигляді функції.

%enditem
}
{\smallskip\smallskip \begin {scriptsize} Затверджено на засіданні кафедри теоретичної кібернетики, протокол \No 4  від   25.01.2024 р.\par\smallskip\smallskip\smallskip\smallskip
{\parbox {6.5 cm} {Зав. кафедри \dotfill Ю.В.~Крак}}\hspace {0.7cm} {\hbox to 6.5 cm { Екзаменатор \dotfill Т.О.~Карнаух}} \end{scriptsize}}
%end
\vspace{0.9cm}
%begin
{{\begin {scriptsize}\par
{ \center  Київський національний університет імені Тараса Шевченка \par}
{\parbox {5 cm}{Освітня програма \hfill}{\parbox {7 cm}{Прикладна математика\hfill}}\parbox {6 cm}{\hfill Семестр 2}}\par
{\parbox {5 cm} {Навчальний предмет \hfill}\parbox {7 cm}{Програмування \hfill}\parbox {6cm}{\hfill}\par}\vspace{-15pt} \end{scriptsize}}{\center Екзаменаційний білет \No \textbf { 35 }\par}}
{%beginitem
1. Алгоритм пошуку в глибину та задачі, які він може вирішувати.%enditem
\par
%beginitem
2. Написати функцію, що в текстовому файлі замінює всі цілі числа на слово NUMBER. Файл дуже великий і його вміст завантажений у пам'ять бути не може. Результат має бути записаний у новий файл.

%enditem
\par
%beginitem
3. Кожен з двох текстових файлів містить послідовність дійсних чисел, записану у форматі одне число на рядок.  Перевірити, чи задають вони одну ту саму послідовність. Кількість чисел у кожному з двох файлів така, що вони не можуть бути завантажені повністю в пам'ять програми. Розв'язання оформити у вигляді функції.

%enditem
}
{\smallskip\smallskip \begin {scriptsize} Затверджено на засіданні кафедри теоретичної кібернетики, протокол \No 4  від   25.01.2024 р.\par\smallskip\smallskip\smallskip\smallskip
{\parbox {6.5 cm} {Зав. кафедри \dotfill Ю.В.~Крак}}\hspace {0.7cm} {\hbox to 6.5 cm { Екзаменатор \dotfill Т.О.~Карнаух}} \end{scriptsize}}
%end
\vspace{0.9cm}
%begin
{{\begin {scriptsize}\par
{ \center  Київський національний університет імені Тараса Шевченка \par}
{\parbox {5 cm}{Освітня програма \hfill}{\parbox {7 cm}{Прикладна математика\hfill}}\parbox {6 cm}{\hfill Семестр 2}}\par
{\parbox {5 cm} {Навчальний предмет \hfill}\parbox {7 cm}{Програмування \hfill}\parbox {6cm}{\hfill}\par}\vspace{-15pt} \end{scriptsize}}{\center Екзаменаційний білет \No \textbf { 36 }\par}}
{%beginitem
1. Дерева. Обходи дерев.%enditem
\par
%beginitem
2. Написати функцію, що в текстовому файлі замінює всі входження фрази \textbf{не зараховано} на слово UPS. Кількість білих символів між частинами фрази може бути довільна (і там можуть бути не тільки пробіли). Результат має бути записаний у новий файл.

%enditem
\par
%beginitem
3. Текстовий файл містить послідовність цілих чисел $a_1$, $a_2$, $\ldots$, $a_t$, записану у форматі одне число на рядок. Знайти сімнадцяте найменше значення, що записане у файлі. Кожне значення, що записане у файлі, враховується стільки разів, скільки входить у файл. Кількість чисел у файлі така, що вони не можуть бути завантажені повністю в пам'ять програми. Розв'язання оформити у вигляді функції.

%enditem
}
{\smallskip\smallskip \begin {scriptsize} Затверджено на засіданні кафедри теоретичної кібернетики, протокол \No 4  від   25.01.2024 р.\par\smallskip\smallskip\smallskip\smallskip
{\parbox {6.5 cm} {Зав. кафедри \dotfill Ю.В.~Крак}}\hspace {0.7cm} {\hbox to 6.5 cm { Екзаменатор \dotfill Т.О.~Карнаух}} \end{scriptsize}}
%end
\newpage
%begin
{{\begin {scriptsize}\par
{ \center  Київський національний університет імені Тараса Шевченка \par}
{\parbox {5 cm}{Освітня програма \hfill}{\parbox {7 cm}{Прикладна математика\hfill}}\parbox {6 cm}{\hfill Семестр 2}}\par
{\parbox {5 cm} {Навчальний предмет \hfill}\parbox {7 cm}{Програмування \hfill}\parbox {6cm}{\hfill}\par}\vspace{-15pt} \end{scriptsize}}{\center Екзаменаційний білет \No \textbf { 37 }\par}}
{%beginitem
1. Контекстні менеджери.%enditem
\par
%beginitem
2. Написати функцію, що за текстовим файлом з статтею визначає скільки разів у ньому містяться слова, що починаються на Київ та Kyiv.  Переноси слів між рядками у файлі відсутні. Файл дуже великий і його вміст завантажений у пам'ять бути не може.

%enditem
\par
%beginitem
3. Текстовий файл містить послідовність цілих чисел $a_1$, $a_2$, $\ldots$, $a_t$, записану у форматі одне число на рядок. Знайти сімнадцяте найменше значення, що записане у файлі. Кожне значення, що записане у файлі, враховується стільки разів, скільки входить у файл. Кількість чисел у файлі така, що вони не можуть бути завантажені повністю в пам'ять програми. Розв'язання оформити у вигляді функції.

%enditem
}
{\smallskip\smallskip \begin {scriptsize} Затверджено на засіданні кафедри теоретичної кібернетики, протокол \No 4  від   25.01.2024 р.\par\smallskip\smallskip\smallskip\smallskip
{\parbox {6.5 cm} {Зав. кафедри \dotfill Ю.В.~Крак}}\hspace {0.7cm} {\hbox to 6.5 cm { Екзаменатор \dotfill Т.О.~Карнаух}} \end{scriptsize}}
%end
\vspace{0.9cm}
%begin
{{\begin {scriptsize}\par
{ \center  Київський національний університет імені Тараса Шевченка \par}
{\parbox {5 cm}{Освітня програма \hfill}{\parbox {7 cm}{Прикладна математика\hfill}}\parbox {6 cm}{\hfill Семестр 2}}\par
{\parbox {5 cm} {Навчальний предмет \hfill}\parbox {7 cm}{Програмування \hfill}\parbox {6cm}{\hfill}\par}\vspace{-15pt} \end{scriptsize}}{\center Екзаменаційний білет \No \textbf { 38 }\par}}
{%beginitem
1. Ітератори.%enditem
\par
%beginitem
2. Написати функцію, що за текстовим файлом, що містить зображення цілих чисел, дійсних чисел та точок площини, розділені білими символами, обчислює кількість входжень у файл зображень точок. Зображення точки має вигляд: відкриваюча дужка, x-координата, кома, y-координата, закриваюча дужка, наприклад, ( 13.5 , 45.1E-1) або ( 14, -3) . Білі символи можуть зустрічатись у довільній кількості. Файл дуже великий і його вміст завантажений у пам'ять бути не може.

%enditem
\par
%beginitem
3. Текстовий файл містить послідовність цілих чисел $a_1$, $a_2$, $\ldots$, $a_t$, записану у форматі одне число на рядок. Знайти тринадцяте найменше значення, що записане у файлі. Кожне значення, що записане у файлі, враховується тільки один раз. Кількість чисел у файлі така, що вони не можуть бути завантажені повністю в пам'ять програми. Розв'язання оформити у вигляді функції.

%enditem
}
{\smallskip\smallskip \begin {scriptsize} Затверджено на засіданні кафедри теоретичної кібернетики, протокол \No 4  від   25.01.2024 р.\par\smallskip\smallskip\smallskip\smallskip
{\parbox {6.5 cm} {Зав. кафедри \dotfill Ю.В.~Крак}}\hspace {0.7cm} {\hbox to 6.5 cm { Екзаменатор \dotfill Т.О.~Карнаух}} \end{scriptsize}}
%end
\vspace{0.9cm}
%begin
{{\begin {scriptsize}\par
{ \center  Київський національний університет імені Тараса Шевченка \par}
{\parbox {5 cm}{Освітня програма \hfill}{\parbox {7 cm}{Прикладна математика\hfill}}\parbox {6 cm}{\hfill Семестр 2}}\par
{\parbox {5 cm} {Навчальний предмет \hfill}\parbox {7 cm}{Програмування \hfill}\parbox {6cm}{\hfill}\par}\vspace{-15pt} \end{scriptsize}}{\center Екзаменаційний білет \No \textbf { 39 }\par}}
{%beginitem
1. Файли формату csv.%enditem
\par
%beginitem
2. Написати функцію, що в текстовому файлі замінює всі зазначення ваги у вигляді \textbf{якесь ціле число кг} на слово WEIGHT (приклади зазначення ваги: 13 кг, 26   кг). Файл дуже великий і його вміст завантажений у пам'ять бути не може. Результат має бути записаний у новий файл.

%enditem
\par
%beginitem
3. Текстовий файл містить послідовність цілих чисел $a_1$, $a_2$, $\ldots$, $a_t$, записану у форматі одне число на рядок. Знайти тринадцяте найбільше значення, що записане у файлі. Кожне значення, що записане у файлі, враховується тільки один раз. Кількість чисел у файлі така, що вони не можуть бути завантажені повністю в пам'ять програми. Розв'язання оформити у вигляді функції.

%enditem
}
{\smallskip\smallskip \begin {scriptsize} Затверджено на засіданні кафедри теоретичної кібернетики, протокол \No 4  від   25.01.2024 р.\par\smallskip\smallskip\smallskip\smallskip
{\parbox {6.5 cm} {Зав. кафедри \dotfill Ю.В.~Крак}}\hspace {0.7cm} {\hbox to 6.5 cm { Екзаменатор \dotfill Т.О.~Карнаух}} \end{scriptsize}}
%end
\newpage
%begin
{{\begin {scriptsize}\par
{ \center  Київський національний університет імені Тараса Шевченка \par}
{\parbox {5 cm}{Освітня програма \hfill}{\parbox {7 cm}{Прикладна математика\hfill}}\parbox {6 cm}{\hfill Семестр 2}}\par
{\parbox {5 cm} {Навчальний предмет \hfill}\parbox {7 cm}{Програмування \hfill}\parbox {6cm}{\hfill}\par}\vspace{-15pt} \end{scriptsize}}{\center Екзаменаційний білет \No \textbf { 40 }\par}}
{%beginitem
1. Обробка файлів: режими відкриття файлів.%enditem
\par
%beginitem
2. Написати функцію, що за текстовим файлом з статтею визначає  скільки разів у ньому містяться слова, що закінчуються на \textbf{ський} або на \textbf{зький} (рід, число та відмінок значення не мають, тобто \textit{міська}, \textit{міському} тощо теж мають рахуватись). Файл дуже великий і його вміст завантажений у пам'ять бути не може.

%enditem
\par
%beginitem
3. Реалізувати операцію різниці множин з повтореннями. Множина з повтореннями задається переліком своїх елементів, записаних за неспаданням у текстовому файлі через білі символи. Результат у такому ж самому форматі (і теж за неспаданням) має бути записаний у текстовий файл. Питома функція має отримувати три шляхи до файлів. Розмір вхідних множин такий, що вони не можуть бути завантажені повністю в пам’ять програми.

%enditem
}
{\smallskip\smallskip \begin {scriptsize} Затверджено на засіданні кафедри теоретичної кібернетики, протокол \No 4  від   25.01.2024 р.\par\smallskip\smallskip\smallskip\smallskip
{\parbox {6.5 cm} {Зав. кафедри \dotfill Ю.В.~Крак}}\hspace {0.7cm} {\hbox to 6.5 cm { Екзаменатор \dotfill Т.О.~Карнаух}} \end{scriptsize}}
%end
\vspace{0.9cm}
%begin
{{\begin {scriptsize}\par
{ \center  Київський національний університет імені Тараса Шевченка \par}
{\parbox {5 cm}{Освітня програма \hfill}{\parbox {7 cm}{Прикладна математика\hfill}}\parbox {6 cm}{\hfill Семестр 2}}\par
{\parbox {5 cm} {Навчальний предмет \hfill}\parbox {7 cm}{Програмування \hfill}\parbox {6cm}{\hfill}\par}\vspace{-15pt} \end{scriptsize}}{\center Екзаменаційний білет \No \textbf { 41 }\par}}
{%beginitem
1. Лінійні зв’язані структури даних.%enditem
\par
%beginitem
2. Написати функцію, що для текстового файлу, що містить послідовність точок площини, обчислює центр мас заданих точок (центр мас -- точка, координати якої є середнім арифметичним відповідних координат). Для кожної точки на окремому рядку наводиться її зображення: відкриваюча дужка, x-координата, кома, y-координата, закриваюча дужка, наприклад, ( 13.5 , 45.1E-1). Білі символи можуть зустрічатись у довільній кількості. Файл дуже великий і його вміст завантажений у пам'ять бути не може.

%enditem
\par
%beginitem
3. Реалізувати операцію злиття множин з повтореннями. Множина з повтореннями задається переліком своїх елементів, записаних за неспаданням у текстовому файлі через білі символи. Результат у такому ж самому форматі (і теж за неспаданням) має бути записаний у текстовий файл. Питома функція має отримувати три шляхи до файлів. Розмір вхідних множин такий, що вони не можуть бути завантажені повністю в пам’ять програми.

%enditem
}
{\smallskip\smallskip \begin {scriptsize} Затверджено на засіданні кафедри теоретичної кібернетики, протокол \No 4  від   25.01.2024 р.\par\smallskip\smallskip\smallskip\smallskip
{\parbox {6.5 cm} {Зав. кафедри \dotfill Ю.В.~Крак}}\hspace {0.7cm} {\hbox to 6.5 cm { Екзаменатор \dotfill Т.О.~Карнаух}} \end{scriptsize}}
%end
\vspace{0.9cm}
%begin
{{\begin {scriptsize}\par
{ \center  Київський національний університет імені Тараса Шевченка \par}
{\parbox {5 cm}{Освітня програма \hfill}{\parbox {7 cm}{Прикладна математика\hfill}}\parbox {6 cm}{\hfill Семестр 2}}\par
{\parbox {5 cm} {Навчальний предмет \hfill}\parbox {7 cm}{Програмування \hfill}\parbox {6cm}{\hfill}\par}\vspace{-15pt} \end{scriptsize}}{\center Екзаменаційний білет \No \textbf { 42 }\par}}
{%beginitem
1. Статичні поля, статичні методи, методи класів.%enditem
\par
%beginitem
2. Написати функцію, що в текстовому файлі замінює всі цілі числа на слово NUMBER. Файл дуже великий і його вміст завантажений у пам'ять бути не може. Результат має бути записаний у новий файл.

%enditem
\par
%beginitem
3. Текстовий файл містить послідовність цілих чисел $a_1$, $a_2$, $\ldots$, $a_t$, записану у форматі одне число на рядок. Знайти сімнадцяте найбільше значення, що записане у файлі. Кожне значення, що записане у файлі, враховується стільки разів, скільки входить у файл. Кількість чисел у файлі така, що вони не можуть бути завантажені повністю в пам'ять програми. Розв'язання оформити у вигляді функції.

%enditem
}
{\smallskip\smallskip \begin {scriptsize} Затверджено на засіданні кафедри теоретичної кібернетики, протокол \No 4  від   25.01.2024 р.\par\smallskip\smallskip\smallskip\smallskip
{\parbox {6.5 cm} {Зав. кафедри \dotfill Ю.В.~Крак}}\hspace {0.7cm} {\hbox to 6.5 cm { Екзаменатор \dotfill Т.О.~Карнаух}} \end{scriptsize}}
%end
\newpage
%begin
{{\begin {scriptsize}\par
{ \center  Київський національний університет імені Тараса Шевченка \par}
{\parbox {5 cm}{Освітня програма \hfill}{\parbox {7 cm}{Прикладна математика\hfill}}\parbox {6 cm}{\hfill Семестр 2}}\par
{\parbox {5 cm} {Навчальний предмет \hfill}\parbox {7 cm}{Програмування \hfill}\parbox {6cm}{\hfill}\par}\vspace{-15pt} \end{scriptsize}}{\center Екзаменаційний білет \No \textbf { 43 }\par}}
{%beginitem
1. Алгоритм пошуку в ширину та задачі, які він може вирішувати.%enditem
\par
%beginitem
2. Написати функцію, що за текстовим файлом з статтею визначає скільки разів у ньому містяться слова, що починаються на Київ та Kyiv.  Переноси слів між рядками у файлі відсутні. Файл дуже великий і його вміст завантажений у пам'ять бути не може.

%enditem
\par
%beginitem
3. Кожен з двох текстових файлів містить послідовність дійсних чисел, записану у форматі одне число на рядок.  Перевірити, чи задають вони одну ту саму послідовність. Кількість чисел у кожному з двох файлів така, що вони не можуть бути завантажені повністю в пам'ять програми. Розв'язання оформити у вигляді функції.

%enditem
}
{\smallskip\smallskip \begin {scriptsize} Затверджено на засіданні кафедри теоретичної кібернетики, протокол \No 4  від   25.01.2024 р.\par\smallskip\smallskip\smallskip\smallskip
{\parbox {6.5 cm} {Зав. кафедри \dotfill Ю.В.~Крак}}\hspace {0.7cm} {\hbox to 6.5 cm { Екзаменатор \dotfill Т.О.~Карнаух}} \end{scriptsize}}
%end
\vspace{0.9cm}
%begin
{{\begin {scriptsize}\par
{ \center  Київський національний університет імені Тараса Шевченка \par}
{\parbox {5 cm}{Освітня програма \hfill}{\parbox {7 cm}{Прикладна математика\hfill}}\parbox {6 cm}{\hfill Семестр 2}}\par
{\parbox {5 cm} {Навчальний предмет \hfill}\parbox {7 cm}{Програмування \hfill}\parbox {6cm}{\hfill}\par}\vspace{-15pt} \end{scriptsize}}{\center Екзаменаційний білет \No \textbf { 44 }\par}}
{%beginitem
1. Емпіричний підхід до визначення складності обчислень.%enditem
\par
%beginitem
2. Написати функцію, що за текстовим файлом, що містить зображення цілих чисел, дійсних чисел та точок площини, розділені білими символами, обчислює кількість входжень у файл зображень точок. Зображення точки має вигляд: відкриваюча дужка, x-координата, кома, y-координата, закриваюча дужка, наприклад, ( 13.5 , 45.1E-1) або ( 14, -3) . Білі символи можуть зустрічатись у довільній кількості. Файл дуже великий і його вміст завантажений у пам'ять бути не може.

%enditem
\par
%beginitem
3. Реалізувати операцію симетричної різниці множин з повтореннями. Множина з повтореннями задається переліком своїх елементів, записаних за неспаданням у текстовому файлі через білі символи. Результат у такому ж самому форматі (і теж за неспаданням) має бути записаний у текстовий файл. Питома функція має отримувати три шляхи до файлів. Розмір вхідних множин такий, що вони не можуть бути завантажені повністю в пам’ять програми.

%enditem
}
{\smallskip\smallskip \begin {scriptsize} Затверджено на засіданні кафедри теоретичної кібернетики, протокол \No 4  від   25.01.2024 р.\par\smallskip\smallskip\smallskip\smallskip
{\parbox {6.5 cm} {Зав. кафедри \dotfill Ю.В.~Крак}}\hspace {0.7cm} {\hbox to 6.5 cm { Екзаменатор \dotfill Т.О.~Карнаух}} \end{scriptsize}}
%end
\vspace{0.9cm}
%begin
{{\begin {scriptsize}\par
{ \center  Київський національний університет імені Тараса Шевченка \par}
{\parbox {5 cm}{Освітня програма \hfill}{\parbox {7 cm}{Прикладна математика\hfill}}\parbox {6 cm}{\hfill Семестр 2}}\par
{\parbox {5 cm} {Навчальний предмет \hfill}\parbox {7 cm}{Програмування \hfill}\parbox {6cm}{\hfill}\par}\vspace{-15pt} \end{scriptsize}}{\center Екзаменаційний білет \No \textbf { 45 }\par}}
{%beginitem
1. Багатовимірні масиви бібліотеки numpy: у чому відмінність моделювання матриці багатовимірним масивом бібліотеки numpy та як список списків.%enditem
\par
%beginitem
2. Написати функцію, що за текстовим файлом з статтею визначає  скільки разів у ньому містяться слова, що закінчуються на \textbf{ський} або на \textbf{зький} (рід, число та відмінок значення не мають, тобто \textit{міська}, \textit{міському} тощо теж мають рахуватись). Файл дуже великий і його вміст завантажений у пам'ять бути не може.

%enditem
\par
%beginitem
3. Реалізувати операцію перетину множин з повтореннями. Множина з повтореннями задається переліком своїх елементів, записаних за неспаданням у текстовому файлі через білі символи. Результат у такому ж самому форматі (і теж за неспаданням) має бути записаний у текстовий файл. Питома функція має отримувати три шляхи до файлів. Розмір вхідних множин такий, що вони не можуть бути завантажені повністю в пам’ять програми.

%enditem
}
{\smallskip\smallskip \begin {scriptsize} Затверджено на засіданні кафедри теоретичної кібернетики, протокол \No 4  від   25.01.2024 р.\par\smallskip\smallskip\smallskip\smallskip
{\parbox {6.5 cm} {Зав. кафедри \dotfill Ю.В.~Крак}}\hspace {0.7cm} {\hbox to 6.5 cm { Екзаменатор \dotfill Т.О.~Карнаух}} \end{scriptsize}}
%end
\newpage
%begin
{{\begin {scriptsize}\par
{ \center  Київський національний університет імені Тараса Шевченка \par}
{\parbox {5 cm}{Освітня програма \hfill}{\parbox {7 cm}{Прикладна математика\hfill}}\parbox {6 cm}{\hfill Семестр 2}}\par
{\parbox {5 cm} {Навчальний предмет \hfill}\parbox {7 cm}{Програмування \hfill}\parbox {6cm}{\hfill}\par}\vspace{-15pt} \end{scriptsize}}{\center Екзаменаційний білет \No \textbf { 46 }\par}}
{%beginitem
1. Математичний підхід до визначення складності обчислень.%enditem
\par
%beginitem
2. Написати функцію, що в текстовому файлі замінює всі зазначення ваги у вигляді \textbf{якесь ціле число кг} на слово WEIGHT (приклади зазначення ваги: 13 кг, 26   кг). Файл дуже великий і його вміст завантажений у пам'ять бути не може. Результат має бути записаний у новий файл.

%enditem
\par
%beginitem
3. Текстовий файл містить послідовність цілих чисел $a_1$, $a_2$, $\ldots$, $a_t$, записану у форматі одне число на рядок. Знайти сімнадцяте найменше значення, що записане у файлі. Кожне значення, що записане у файлі, враховується стільки разів, скільки входить у файл. Кількість чисел у файлі така, що вони не можуть бути завантажені повністю в пам'ять програми. Розв'язання оформити у вигляді функції.

%enditem
}
{\smallskip\smallskip \begin {scriptsize} Затверджено на засіданні кафедри теоретичної кібернетики, протокол \No 4  від   25.01.2024 р.\par\smallskip\smallskip\smallskip\smallskip
{\parbox {6.5 cm} {Зав. кафедри \dotfill Ю.В.~Крак}}\hspace {0.7cm} {\hbox to 6.5 cm { Екзаменатор \dotfill Т.О.~Карнаух}} \end{scriptsize}}
%end
\vspace{0.9cm}
%begin
{{\begin {scriptsize}\par
{ \center  Київський національний університет імені Тараса Шевченка \par}
{\parbox {5 cm}{Освітня програма \hfill}{\parbox {7 cm}{Прикладна математика\hfill}}\parbox {6 cm}{\hfill Семестр 2}}\par
{\parbox {5 cm} {Навчальний предмет \hfill}\parbox {7 cm}{Програмування \hfill}\parbox {6cm}{\hfill}\par}\vspace{-15pt} \end{scriptsize}}{\center Екзаменаційний білет \No \textbf { 47 }\par}}
{%beginitem
1. Множини.%enditem
\par
%beginitem
2. Написати функцію, що в текстовому файлі замінює всі цілі числа на слово NUMBER. Файл дуже великий і його вміст завантажений у пам'ять бути не може. Результат має бути записаний у новий файл.

%enditem
\par
%beginitem
3. Кожен з двох текстових файлів містить послідовність дійсних чисел, записану у форматі одне число на рядок.  Перевірити, чи задають вони одну ту саму послідовність. Кількість чисел у кожному з двох файлів така, що вони не можуть бути завантажені повністю в пам'ять програми. Розв'язання оформити у вигляді функції.

%enditem
}
{\smallskip\smallskip \begin {scriptsize} Затверджено на засіданні кафедри теоретичної кібернетики, протокол \No 4  від   25.01.2024 р.\par\smallskip\smallskip\smallskip\smallskip
{\parbox {6.5 cm} {Зав. кафедри \dotfill Ю.В.~Крак}}\hspace {0.7cm} {\hbox to 6.5 cm { Екзаменатор \dotfill Т.О.~Карнаух}} \end{scriptsize}}
%end
\vspace{0.9cm}
%begin
{{\begin {scriptsize}\par
{ \center  Київський національний університет імені Тараса Шевченка \par}
{\parbox {5 cm}{Освітня програма \hfill}{\parbox {7 cm}{Прикладна математика\hfill}}\parbox {6 cm}{\hfill Семестр 2}}\par
{\parbox {5 cm} {Навчальний предмет \hfill}\parbox {7 cm}{Програмування \hfill}\parbox {6cm}{\hfill}\par}\vspace{-15pt} \end{scriptsize}}{\center Екзаменаційний білет \No \textbf { 48 }\par}}
{%beginitem
1. Файли формату json.%enditem
\par
%beginitem
2. Написати функцію, що для текстового файлу, що містить послідовність точок площини, обчислює центр мас заданих точок (центр мас -- точка, координати якої є середнім арифметичним відповідних координат). Для кожної точки на окремому рядку наводиться її зображення: відкриваюча дужка, x-координата, кома, y-координата, закриваюча дужка, наприклад, ( 13.5 , 45.1E-1). Білі символи можуть зустрічатись у довільній кількості. Файл дуже великий і його вміст завантажений у пам'ять бути не може.

%enditem
\par
%beginitem
3. Текстовий файл містить послідовність цілих чисел $a_1$, $a_2$, $\ldots$, $a_t$, записану у форматі одне число на рядок. Знайти сімнадцяте найбільше значення, що записане у файлі. Кожне значення, що записане у файлі, враховується стільки разів, скільки входить у файл. Кількість чисел у файлі така, що вони не можуть бути завантажені повністю в пам'ять програми. Розв'язання оформити у вигляді функції.

%enditem
}
{\smallskip\smallskip \begin {scriptsize} Затверджено на засіданні кафедри теоретичної кібернетики, протокол \No 4  від   25.01.2024 р.\par\smallskip\smallskip\smallskip\smallskip
{\parbox {6.5 cm} {Зав. кафедри \dotfill Ю.В.~Крак}}\hspace {0.7cm} {\hbox to 6.5 cm { Екзаменатор \dotfill Т.О.~Карнаух}} \end{scriptsize}}
%end
\newpage
\end{document}
